\documentclass[11pt]{article}
\usepackage[margin=1in]{geometry}
\usepackage{amsmath,hyperref}
\title{When recovery consumes a support handoff's workcell}
\date{}
\begin{document}
\maketitle
The connected warehouse workload in \href{https://github.com/eternal-november/eos-research-programmes/pull/356}{PR 356} makes the root's relief depend on an authorization and a downstream load transfer. Here an independently owed inspection, not longer transfer dwell, can make the necessary transfer late. The inspection's recovery changes which crew controls the timing of an otherwise unrelated handoff. This is a second authored workload on the same event platform, not a new independently chosen transfer. The scope is a finite civil warehouse model, not safety certification for equipment.

\paragraph{Plant, parties, and reached cuts.}
The warehouse owes one relief of load 0 by tick 5 and continuous support for every unsettled load. In this workload the rack owner is also owed an inspection of supported load 3 by tick 4. Before either incident, the two commitments are logged at tick $-1$. Root authority is a named maintenance order; crew 3 has the scheduled inspection and retry permissions, neither confers root authority. The original placement and rack settlement events produce the same short cut (no live load 1) and long cut (live load 1), both with loads 0 and 3 still supported. In the long cut crew 1 cannot accept load 0 until its own load is accepted at support 2. Support transfers are one tick, nonpreemptive, retain the old support up to acceptance, and reserve both endpoints. No old commitment is dropped and no new support is added.

Only the 1-to-2 transfer and rack inspection use one physical workcell lane. Inspection occupies it during one tick, without moving load 3. The terminal controller sees lane availability at its installed local admission guard together with the same adjacent stocks and authorization. The root receiver and channel are the pinned code. Local interlock grants on actual resources, not knowledge of the global scenario. For normal service crew 3 attempts inspection at tick 1 and releases the lane at tick 2. A failed attempt triggers its preauthorized retry in the sole permitted contingency window $[3,4)$. Retry failure and discretion to shift that window are outside this envelope. Due acceptances and releases precede later starts. On retry tick the interlock admits the fixed-slot inspection before considering any new handoff action. The additional dependency is not a scheduler giving the dispatcher private knowledge.

\paragraph{Four matched arms.}
Change A is the primary one-tick first-hop queue delay at tick 0. Change B is a transient jam of the inspection head at tick 1 before inspection can start; only crew 3 sees this failed start. Crew 3's ordinary procedure then attempts the owed inspection at 3. Crew 2 sees only the lane unavailable at that cut. Both events and the work schedule are fixed independently of success. No message declares the true arm. Unjammed operation uses $[1,2)$; recovery uses $[3,4)$. The same crew-2 and crew-3 policies run in every arm.
\[
\begin{array}{c|rrrr}
&\text{nominal}&A&B&A+B\\\hline
\text{long, incumbent root finish}&4&5&4&\text{miss}\\
\text{short, incumbent root finish}&2&3&2&3\\
\text{long, early issue, onset first}&4&4&3&4\\
\text{long, early issue, delivery first}&4&4&3&3\\
\text{inspection finish}&2&2&4&4
\end{array}
\]
In the $A+B$ long arm, the first permit reaches 1 at tick 2 and the forwarded permit reaches 2 at tick 3. The retry already holds the shared lane then. This is the first new downstream edge: recovery at crew 3 removes a start at the 1-to-2 interlock. Once the lane releases at 4, transfer 1-to-2 starts, accepts at 5, and finally makes support 1 free. The second downstream edge is the support release consumed by the root's unchanged admission rule; the root now has no time left. Load 0 remains safely supported but owed relief fails. The short cut has the same faults and same retry but the settled load 1 removes the dependency on the lane.

An intervention on just this coupling keeps both faults, the same failed attempt, the same retry timing and the same acceptances for inspection. Giving inspection a physically independent unit lane restores root relief at tick 5. This injects an extra lane for a mediation control only; it is not the allowed repair. Unlike turning off the jam, it preserves the initiating events and recovery duty. The comparison to the earlier workload is substantive: load 3, previously an untouched support, now has a job that contends with a transfer. Both workflows run on the same pinned event engine, and the second workload preserves the original root service and safety test while adding the explicit inspection promise.

\paragraph{Repair, lower bound, and cost.}
The installed repair preissues the already authorized order at $-1$ and does not permit any action before tick 0. Under onset-first ordering A moves the packet due at 0 to 1; under delivery-first ordering it is already received at 0 and cannot be retimed. Crew 1 forwards on actual receipt, and crew 2 starts only after local authorization and an available lane. The table and exhaustive enumeration give the deterministic policy on all admitted combinations at deadline 5. Crew 3 still inspects and crew 1's code is unchanged; its authority and support timing can change. There is one earlier slot of communication commitment, no extra packet, lane, support, service refusal or cancelled obligation. Preauthorization is a material assumption.

Let $a$ be crew 1's earliest permit receipt and $I$ the unique inspection slot. For a long cut, every lawful schedule must start the terminal move at $s\ge a+1$, with $s\notin I$, and root relief cannot finish before $s+2$. The installed greedy rule attains this bound whenever the deadline permits its final start. In the narrower class with no issue before 0, unchanged channel minima, no bypass, unchanged inspection window, and the same interlock, $a\ge2$ under A and $s\ge4$ under $A+B$. Even a clairvoyant scheduler in that class has earliest relief 6. At root deadline 6 the combined arm succeeds; at deadline 4 hop-only also fails. This is a sharp event boundary, not uniform tolerance of other retry slots or prolonged dwell. For example, permitting a legal shift of inspection to tick 2 can remove this obstruction, but changes the stated contingency law. If only the repair's communication lead changes, all original resource constraints remain.

\paragraph{Proof comparison and remaining transfer scope.}
For an authored check, instantiate the frozen \href{https://github.com/eternal-november/eos-research-programmes/blob/e1348964f73704212d2d8b48d5982a5ed600a2ff/projects/contextual-composition/feedback-and-concurrency/observable-covers/Residual_Profiles_and_Admission_Reflection.tex}{residual-profile theorem} on the original one-tick-dwell subcarrier. For each reached local phase the cut key is (old stock workload, one-tick terminal dwell, deadline, tick). Map witnesses by that key to nominal primary receiver witnesses; keep source process identity as provenance, not a read. Both cuts and deadlines 4 through 6 are compared. The recovered combined arm with early issue and onset first has exactly the same named permit and visible support traces as nominal primary. The old retry persists in the actual event history but has no effect on those reads because the transfer is admitted before the retry acquires the lane. Retain the exact prior phase extractors and tests; replay compares all raw factor arguments including incompatible cross-key tuples. Unary reads distinguish payloads, and the cross-key tests distinguish the rest; hence target residual classes are (key, payload) and have a map preimage. The phase extraction separately checks actual reachability and next-tick acceptance; observables alone do not establish these.

The same keyed strategy comparing the combined incumbent to primary hop-only keeps authority receipts equal but fails the physical factor. In the long tick-4 phase primary support 1 is free and has its terminal acceptance, whereas the retry workload still holds load 1. A certificate about authorization alone, or a transported guarantee under a map ignoring this discrepancy, fails. The cached modular direct proof requires the same channel, lane, support, and next-tick acceptance facts; it uses those checked inputs without adding the class coverage proof. The contextual route adds the maps and raw tests; no savings follow.

The new workload supplies changed physical dependence with an authored proof transport and rejection, not a new held-out result. The independently selected positive and rejected bridge cases in PR 356 remain independent and unchanged, and cover the same installed receiver on the original workload. If integrated review requires the independent test itself to change physical dependence on the second workload, commission a non-author after freezing this extended envelope, rather than silently counting this comparison. The admissible frozen interface must retain named authority, schedule and physical support observations; a reviewer can independently choose a different source of workcell contention without new receiver evidence. It must establish actual projected histories and reflection or exhibit their precise obstruction, with a fair direct proof on the same phases. No selection is made here.

\paragraph{Evidence.}
The case register records inventory mechanisms and reopen conditions. The replay compiles the exact pinned source, changes only the documented construction and admission phases, and compares 96 runs against a separate feasible-tick enumerator. Independent trace assertions check the lane ledger, completion of the inherited inspection, and both faults in the extra-lane control. The receiver check compares 48 phase witnesses per component, 4,608 binary tests and 110,592 raw start triples per direction. The proof comparison does not replace non-author review or assert full contract acceptance.
\end{document}
